\bookchapter{Capítulo 54: Projeto 1 — Meu controle de gastos (versão local)}

\chapterintro{Bem-vindo ao seu primeiro 'projeto sério'. Este controle de gastos, embora não se conecte à internet, tem tudo o que precisa: registrar um lançamento, ver o extrato, somar por mês, excluir e editar. Ao terminá-lo, você terá percorrido o fluxo completo de 'dados lógica interface'. Não tenha pressa de copiar código — pense primeiro, aja depois; esse é o nosso princípio de sempre.}

\booksection{Pensando nos requisitos: faça primeiro cinco perguntas}

Antes de pôr a mão na massa, organize a ideia com os 'cinco elementos de requisito':

\begin{center}
\begin{tabularx}{\linewidth}{|X|X|}
\hline
\textbf{Elemento} & \textbf{Resposta} \tabularnewline
\hline
Who (para quem é) & Para você mesmo; família e amigos também valem \tabularnewline
What (o que faz) & Registrar receitas e despesas, ver o extrato, resumir por mês, estatísticas por categoria \tabularnewline
Why (por que faz) & Praticar + descobrir para onde o dinheiro está indo \tabularnewline
How (como faz) & Kotlin + Compose + Room + MVVM, 100\% local, sem back-end \tabularnewline
Constraint (restrições) & Concluir em uma semana; suportar no mínimo o Android 8; mesmo sem publicar, já escrever no padrão de publicação \tabularnewline
\hline
\end{tabularx}
\end{center}

\begin{center}\small\color{gray}Os cinco elementos de requisito do controle de gastos\end{center}

Depois, divida em passos pequenos em que 'cada etapa já serve para algo': (1) criar a camada de dados com Room (conseguir salvar um lançamento); (2) tela de lista (conseguir ver o extrato); (3) tela de entrada (conseguir adicionar); (4) excluir e editar; (5) cartão de resumo mensal. Ao fim de cada passo, o App já é uma 'coisa utilizável'.

\booksection{Camada de dados: dando aos 'lançamentos' um lar}

Primeiro defina como é um lançamento — campos suficientes, mas sem gula:

\begin{lstlisting}[language=Kotlin]
@Entity(tableName = "records")
data class Record(
    @PrimaryKey(autoGenerate = true) val id: Long = 0,
    val type: String,        // "income" receita | "expense" despesa
    val amount: Double,      // valor (em yuan)
    val category: String,    // categoria: alimentação, transporte, salário…
    val note: String = "",   // observação (opcional)
    val date: String         // "2026-08-15"
)
\end{lstlisting}

No Dao, acrescente duas consultas 'de estatística' — elas são a alma do controle de gastos:

\begin{lstlisting}[language=Kotlin]
@Dao
interface RecordDao {
    @Query("SELECT * FROM records ORDER BY date DESC, id DESC")
    fun getAll(): Flow<List<Record>>

    @Insert
    suspend fun insert(record: Record)

    @Delete
    suspend fun delete(record: Record)

    @Query("SELECT COALESCE(SUM(amount),0) FROM records WHERE type='income' AND date LIKE :month || '%'")
    suspend fun monthIncome(month: String): Double   // month no formato "2026-08"

    @Query("SELECT COALESCE(SUM(amount),0) FROM records WHERE type='expense' AND date LIKE :month || '%'")
    suspend fun monthExpense(month: String): Double
}
\end{lstlisting}

\begin{tcbinfo}{Dica: por que usar Double para o valor?}
Apps financeiros rigorosos usam BigDecimal ou 'Int em centavos' para evitar erros de ponto flutuante (o velho problema de 0,1+0,2$\neq$0,3). Para um projeto de estudo, Double basta; se um dia você fizer um produto financeiro de verdade, lembre de trocar para centavos como unidade — é uma armadilha que inúmeros antecessores pisaram.

\end{tcbinfo}

\booksection{MVVM: interface e lógica 'morando separadas'}

Deixar a interface operar o banco direto? Funciona, mas o código da tela fica enorme e desagradável de ler. A arquitetura MVVM reparte as tarefas: a View (interface) só desenha, o ViewModel (centro da lógica) cuida dos dados e o Model (Room) cuida do armazenamento.

View (interface Compose)

ViewModel (estado e lógica)

Model (camada de dados Room)

\begin{center}\small\color{gray}Figura 5: A estrutura de três camadas do MVVM — os dados fluem em um único sentido e cada camada não invade a outra\end{center}

\begin{lstlisting}[language=Kotlin]
class RecordViewModel(private val dao: RecordDao) : ViewModel() {

    // Converte o Flow em estado observável pela tela
    val records = dao.getAll()
        .stateIn(viewModelScope, SharingStarted.WhileSubscribed(), emptyList())

    // Ações do usuário: chamar funções suspend
    fun addRecord(type: String, amount: Double, category: String, note: String, date: String) {
        viewModelScope.launch {
            dao.insert(Record(type = type, amount = amount,
                category = category, note = note, date = date))
        }
    }

    fun deleteRecord(record: Record) {
        viewModelScope.launch { dao.delete(record) }
    }
}
\end{lstlisting}

Na interface (Compose), use assim: val vm: RecordViewModel = viewModel() e depois val list by vm.records.collectAsState() — assim que os dados mudam, a lista se atualiza sozinha e você não precisa fazer nada. Esse é o poder da união de MVVM + Flow + Compose.

\booksection{Interface: tela de lista + tela de entrada}

Na tela de lista, use LazyColumn (lista de carregamento preguiçoso) — nem com muitos dados ela engasga:

\begin{lstlisting}[language=Kotlin]
@Composable
fun RecordListScreen(vm: RecordViewModel) {
    val records by vm.records.collectAsState()

    Column {
        // card de resumo no topo (receitas - despesas = saldo)
        SummaryCard(records)

        // lista de lançamentos
        LazyColumn {
            items(records, key = { it.id }) { record ->
                RecordRow(record) { vm.deleteRecord(record) }
            }
        }

        // botão flutuante "nova entrada"  abre o diálogo de entrada
        FloatingActionButton(onClick = { showAddDialog = true }) {
            Text("+")
        }
    }
}

@Composable
fun RecordRow(record: Record, onDelete: () -> Unit) {
    Row(Modifier.fillMaxWidth().padding(12.dp)) {
        Text(record.category, fontWeight = FontWeight.Bold)
        Spacer(Modifier.weight(1f))
        Text(if (record.type == "income") "+${record.amount}" else "-${record.amount}",
             color = if (record.type == "income") Color(0xFF27AE60) else Color(0xFFC0392B))
    }
}
\end{lstlisting}

Para a entrada, o mais prático é um AlertDialog (caixa de diálogo): alguns campos de texto + alternância receita/despesa + seletor de categorias; ao tocar em confirmar, chame vm.addRecord(...).

\begin{tcbwarn}{Atenção: não seja preguiçoso na validação de entradas}
O campo de valor precisa de validação: bloqueie entradas inválidas como vazio, número negativo ou 'abc' (use toDoubleOrNull() para verificar). Uma das causas mais comuns de App que cai é 'tratar a entrada do usuário como um anjinho'. Lembre-se: tudo o que o usuário digita deve ser examinado como se fosse inimigo.

\end{tcbwarn}

\booksection{Editar e excluir: deixando os dados 'entrar e sair'}

Um controle de gastos que só adiciona e nunca exclui morre em três dias de uso. Complete com duas operações:

\begin{lstlisting}[language=Kotlin]
// Adicione mais dois ao Dao:
@Query("UPDATE records SET amount=:amount, category=:category, note=:note WHERE id=:id")
suspend fun update(id: Long, amount: Double, category: String, note: String)

@Query("DELETE FROM records WHERE id=:id")
suspend fun deleteById(id: Long)

// Na tela: toque longo no item abre o menu "editar/excluir" (no Compose use DropdownMenu ou AlertDialog)
// É melhor mostrar confirmação antes de excluir — exclusão acidental é o que usuários de apps de despesas mais odeiam
\end{lstlisting}

Duas lições sobre detalhes de interação: excluir exige confirmação dupla (para evitar dedos escorregadios) e editar deve preencher previamente o valor original (corrigir uma letra é mais fácil que digitar tudo de novo). Não menospreze esses detalhes — é deles que nasce a sensação de produto.

\booksection{Aceitação: três níveis de checagem + dois desafios avançados}

\booksection{Aceitação: três níveis de checagem + dois desafios avançados}

Projeto terminado, faça a aceitação em três níveis 'função limite arquitetura': dá para registrar? dá para excluir? ao reabrir o App, os dados continuam lá (função); digitar 0 ou negativo no valor derruba o App? girar a tela perde dados? (limites); interface e dados estão misturados? campos estão hardcoded? (arquitetura).

Depois de passar na aceitação da versão básica, sirva-se de dois 'pratos avançados':

\begin{enumerate}
  \item Filtro por mês: adicione no topo um seletor de mês (por exemplo, 2026-08 v) e a lista mostra só o extrato do mês;
  \item Gráfico de pizza por categoria: use o Canvas do Compose ou uma biblioteca de gráficos leve para desenhar 'para onde foi o dinheiro' numa pizza — com isso pronto, o controle de gastos ganha cara de 'produto'.
\end{enumerate}

\begin{tcbpractice}{Mãos à obra: termine-o nesta semana}
Trabalhe na ordem 'camada de dados ViewModel tela de lista tela de entrada cartão de resumo', um passo por dia. Se travar, jogue a mensagem de erro tal como está no buscador ou na IA (lembra da Parte 7, 'Prática de Programação com IA'? Ao descrever o problema, deixe claros 'comportamento esperado, comportamento real e passos para reproduzir'). Depois de pronto, use-o para registrar seus gastos por uma semana — usar um App que você mesmo fez dá uma sensação de conquista que empurra você a seguir em frente.

\end{tcbpractice}

\begin{tcbsummary}{Resumo do capítulo}
\begin{itemize}
  \item Cinco elementos de requisito dividir em passos pequenos em que 'cada etapa já serve' e só então agir
  \item Camada de dados: entidade Record + Dao (com SQL de estatística mensal) + Database
  \item Divisão MVVM: a View só desenha, o ViewModel cuida da lógica, o Model do armazenamento
  \item Flow + stateIn + collectAsState: os dados mudam e a interface se atualiza sozinha
  \item LazyColumn para a lista de carregamento preguiçoso; AlertDialog para a entrada de dados
  \item Aceitação em três níveis (função/limite/arquitetura) + desafios avançados (filtro mensal, pizza por categoria)
\end{itemize}

\end{tcbsummary}
